\chapter{Gardnerella and Bacterial Vaginosis}
\index{Gardnerella and Bacterial Vaginosis}

\section*{Definition and Overview}
Gardnerella vaginalis is a bacterium that normally lives in small numbers in the vagina, but when the balance of vaginal bacteria is disturbed, it can overgrow and cause bacterial vaginosis (BV). Bacterial vaginosis is the most common cause of abnormal vaginal discharge, producing a thin, greyish-white discharge with a characteristic fishy smell, especially after sex. It is not a sexually transmitted infection in the strict sense, but sexual activity increases the risk. Bacterial vaginosis is common and easily treated with antibiotics, but it recurs in a large proportion of women, and it is associated with pregnancy complications and increased risk of some infections if untreated. This chapter explains the condition, its treatment, and how to reduce recurrence.

\section*{Epidemiology}
Bacterial vaginosis affects around 20--30% of women of reproductive age, making it the most common cause of abnormal vaginal discharge. Rates are higher in sexually active women, in women who use vaginal douching, and with certain hormonal and lifestyle factors. Around half of women with BV have no symptoms. Recurrence is common: around half of treated women have a recurrence within a year. BV is associated with an increased risk of pelvic inflammatory disease, of infection after gynaecological procedures, and of preterm birth in pregnancy, which is why it is taken seriously.

\section*{Aetiology and Pathophysiology}
\begin{itemize}
    \item \textbf{The vaginal microbiome:} a healthy vagina is dominated by lactobacilli, which produce acid that keeps harmful bacteria in check
    \item \textbf{The imbalance:} in BV, lactobacilli decline and Gardnerella and other bacteria overgrow, changing the vaginal acidity and producing amines that cause the fishy odour
    \item \textbf{Risk factors:} new or multiple sexual partners, douching, smoking, hormonal changes, and intrauterine device use
    \item \textbf{Not a classic STI:} BV is not caused by a single sexually transmitted organism, but sex changes the vaginal environment and triggers it
    \item \textbf{Transmission and partners:} treating male partners does not reduce recurrence, and it is not considered a sexually transmitted infection requiring partner treatment
    \item \textbf{Effects:} the altered flora increases susceptibility to chlamydia, gonorrhoea, HIV, and pelvic inflammatory disease
    \item \textbf{Pregnancy effects:} BV is associated with preterm birth and other pregnancy complications
    \item \textbf{Mechanism of symptoms:} the amines and altered pH produce the discharge and odour; inflammation is usually mild, which is why itching and soreness are less prominent than in thrush
\end{itemize}

\section*{Clinical Features}
\begin{itemize}
    \item A thin, greyish-white or milky vaginal discharge
    \item A fishy odour, most noticeable after intercourse and during menstruation
    \item No itching or soreness in most cases, which distinguishes it from thrush
    \item In some women, no symptoms at all
    \item The discharge may be heavier at certain times of the cycle
    \item In pregnancy, it may be detected at antenatal screening
    \item Itching, burning, or pain on urination in some women
\end{itemize}

\section*{Differential Diagnosis}
Vaginal discharge has several causes.
\begin{itemize}
    \item Thrush (candidiasis), with thick white discharge, itching, and soreness
    \item Trichomoniasis, a sexually transmitted infection with frothy, yellow-green, smelly discharge and itching
    \item Chlamydia and gonorrhoea, with discharge and possible pelvic pain
    \item Normal physiological discharge, which varies through the cycle
    \item Desquamative inflammatory vaginitis and other rare causes
    \item A swab and examination distinguish these reliably
\end{itemize}

\section*{When to Seek Medical Care}
Women with a new or unusual discharge, a fishy odour, or discharge with pelvic pain should be assessed. Pregnant women with symptoms should seek review, and women having gynaecological procedures may be screened.

\textbf{Contact your local emergency services for:}
\begin{itemize}
    \item Severe pelvic or abdominal pain with fever, which may indicate pelvic inflammatory disease
    \item Bleeding in pregnancy with pain
    \item Signs of sepsis after a procedure
    \item Any emergency symptoms
\end{itemize}

\section*{Diagnosis and Investigations}
\begin{itemize}
    \item \textbf{History and examination:} the discharge's appearance, odour, and associated symptoms
    \item \textbf{Speculum examination:} assessment of the discharge and vaginal walls
    \item \textbf{Microscopy:} a wet mount shows clue cells (vaginal cells coated with bacteria), the diagnostic feature
    \item \textbf{Vaginal pH testing:} a pH above 4.5 supports BV
    \item \textbf{Whiff test:} adding potassium hydroxide releases the fishy amine odour
    \item \textbf{Swabs for other infections:} chlamydia and gonorrhoea screening when risk factors are present
    \item \textbf{Testing in pregnancy:} screening is considered in some circumstances because of the association with preterm birth
\end{itemize}

\section*{Management}
\begin{itemize}
    \item \textbf{Antibiotics:} metronidazole tablets or vaginal gel, or clindamycin cream, given as a short course
    \item \textbf{Treating symptoms:} the discharge and odour resolve within days of treatment
    \item \textbf{Recurrence:} repeat courses are common; a longer course or maintenance treatment may be offered for frequent recurrence
    \item \textbf{Avoid alcohol with metronidazole:} it causes severe nausea and flushing
    \item \textbf{In pregnancy:} treatment is given for symptomatic BV and in some cases to reduce preterm birth risk, under obstetric advice
    \item \textbf{Before procedures:} treating BV before surgery or intrauterine device insertion reduces complications
    \item \textbf{Partner treatment:} not required, as it does not prevent recurrence
    \item \textbf{Preventive measures:} stop douching, limit the number of sexual partners, and consider condom use
\end{itemize}

\section*{Complications}
\begin{itemize}
    \item Recurrence, which affects around half of women within a year
    \item Increased risk of pelvic inflammatory disease and of infection after procedures
    \item Increased susceptibility to STIs, including HIV
    \item Pregnancy complications, including preterm birth and low birth weight
    \item The discomfort and embarrassment of persistent discharge and odour
\end{itemize}

\section*{Prognosis and Outcomes}
Bacterial vaginosis is easily treated, and symptoms resolve in almost all women with a course of antibiotics. The main challenge is recurrence, which affects a large proportion of women and may require repeated or longer treatment. BV does not cause lasting harm to most women, and the pregnancy and infection risks are managed with appropriate treatment and screening. With good hygiene practices and the avoidance of douching, many women reduce the frequency of episodes.

\section*{Prevention and Self-Care}
\begin{itemize}
    \item Do not douche or wash inside the vagina, which disturbs the natural balance
    \item Use plain water or gentle, unscented products for external washing only
    \item Use condoms, which reduce the risk of changes to the vaginal flora
    \item Limit the number of sexual partners
    \item Avoid smoking
    \item Complete the full course of antibiotics, and avoid alcohol with metronidazole
    \item See a doctor for recurrence rather than self-treating
    \item Attend antenatal care, where BV can be detected and treated
\end{itemize}

\section*{Sources and Further Reading}
The following reputable sources provide further information for healthcare students and patients seeking reliable, current advice about Gardnerella and bacterial vaginosis.
\begin{itemize}
    \item Healthdirect Australia. \textit{Bacterial vaginosis.} https://www.healthdirect.gov.au/bacterial-vaginosis
    \item The Royal Women's Hospital. \textit{Bacterial vaginosis.} https://www.thewomens.org.au/
    \item Jean Hailes for Women's Health. \textit{Bacterial vaginosis.} https://www.jeanhailes.org.au/
    \item NHS. \textit{Bacterial vaginosis.} https://www.nhs.uk/conditions/bacterial-vaginosis/
    \item Mayo Clinic. \textit{Bacterial vaginosis.} https://www.mayoclinic.org/
    \item MedlinePlus. \textit{Bacterial vaginosis.} https://medlineplus.gov/
\end{itemize}
